\section{Rational Scholl realization and Hodge numbers}\label{sec:scholl}
Write $\zeta=\zeta_K$. In this section $\pi$ also denotes the pullback
of the Legendre family to the displayed open base.
For any prime $\ell$, put
\[
 H_\ell=H^1_!(T_{\overline\Q},\Sym^4 R^1\pi_*\Q_\ell),\qquad
 W_{4,\ell}^{\pm}=(e_{\nt}H_\ell)^{I=\pm1},
\]
using the same literal deck and inversion maps as in
Section~\ref{sec:descent}. At $\ell=2$ these are the already constructed
$W_4^\pm$, without a change of sign or a Tate twist.

\begin{theorem}[Rational Scholl realization]\label{thm:scholl-realization}
There are a smooth projective, possibly disconnected fivefold $X/\Q$
and orthogonal idempotents
$p_\pm\in\CH^5(X\times_\Q X)\otimes\Q$ such that for every prime $\ell$
\[
 p_\pm H^j_{\et}(X_{\overline\Q},\Q_\ell)=0\quad(j\ne5),\qquad
 p_\pm H^5_{\et}(X_{\overline\Q},\Q_\ell)\simeq W_{4,\ell}^\pm.
\]
The corresponding Betti realization is compatible with Betti--etale
comparison and has $h^{5,0}=h^{0,5}=2$, with all other Hodge numbers zero.
Thus $(X,p_\pm,0)$ are rational Chow motives realizing the normalized
descents. The construction does not require irreducibility.
\end{theorem}
The following construction proves the theorem. Its use of Scholl's
standard projector theorem is separated from the explicit cover, rational
descent and sign computations for the present family.

\subsection{A semistable auxiliary surface}
Use the finite etale double cover
\[
S=\PP^1_s\setminus\{0,\infty,\mu_6\}\longrightarrow T,
\qquad z=s^2,\qquad \lambda=s^6.
\]
On $S$ the elliptic family is $E_s:y^2=x(x-1)(x-s^6)$. Its Weierstrass
invariants are
\[
\Delta=16s^{12}(1-s^6)^2,\qquad c_4=16(1-s^6+s^{12}).
\]
In characteristic zero, these give an $I_{12}$ fibre at $s=0$ and an $I_2$
fibre at each of the six roots of $s^6=1$. At infinity put
$t=1/s$, $X=t^6x$, $Y=t^9y$; the equation becomes
$Y^2=X(X-1)(X-t^6)$, giving another $I_{12}$ fibre.
The minimal regular model is therefore semistable, with no other singular
fibres. Its nonconstant $j$-invariant makes the family strictly nonconstant.

Now take the \emph{full frame torsor}
\[
S'=\operatorname{Isom}_S((\Z/12\Z)^2,E[12])\longrightarrow S.
\]
It includes all frames, not just one fixed Weil-pairing component and not
frames modulo their sign. It is a finite etale torsor under
$G=\GL_2(\Z/12\Z)$, defined over $\Q$, although it can be disconnected.
Let $\overline S'$ be its smooth projective normalization over $\PP^1_s$,
and let $\overline E'\to\overline S'$ be the minimal regular elliptic model.
Write $\pi':\overline E'|_{S'}\to S'$ for its smooth restriction, and
$\omega'$ for the Hodge line bundle on $\overline S'$.

\textbf{Why level 12 meets every component.} Around a geometric $I_m$ cusp,
the local monodromy on a full level-12 frame is conjugate to
\[
\begin{pmatrix}1&m\\0&1\end{pmatrix}\pmod {12}.
\]
Its order is $e=12/\gcd(12,m)$, the geometric ramification index of $S'$.
After base change and minimal regularization the fibre is $I_{me}$.
Here $m=12$ or $2$, so $me=12$ in both cases. In the Tate description,
the points represented by $q^{a/12}\zeta_{12}^b$ meet all twelve components
of an $I_{12}$ fibre. The tautological frame makes these points sections.
Thus the constant group $V=(\Z/12\Z)^2$ extends into the smooth locus and
meets every geometric fibre component. Any splitting needed in this local
description can be checked after a finite extension; the torsor itself is rational.

These facts verify the two geometric hypotheses of Scholl \cite[Section~2.2]{SchollVanishing}.
On each geometric component one can first enlarge its number field of
constants, apply that result, and then retain the entire finite disjoint union.
No connected-component choice is needed for the rational construction.

Form $(\overline E')^4_{\overline S'}$ and apply Scholl's canonical resolution;
call the result $X$. It has dimension five. The resolution blows up the
intrinsic strata determined by how many elliptic coordinates are nodes
\cite[Section~3.0]{SchollMotives}. The nodal local models suffice \cite[Section~2.10]{SchollVanishing}. Consequently every
family isomorphism considered below preserves its centers and lifts to $X$.


\subsection{Literal automorphisms and the sign relation}
On the family over $S$, define
\[
\begin{aligned}
a(s,x,y)&=(\zeta^2s,x,y),& d(s,x,y)&=(-s,x,y),\\
b(s,x,y)&=(s^{-1},x/s^6,y/s^9),& c(s,x,y)&=(s,x,-y).
\end{aligned}
\]
Use the same maps diagonally on four fibre coordinates. The order-three
choice $a(s)=\zeta^2s$ covers the required $z\mapsto\zeta z$.
For $b$, the chosen square root is exactly $u=s^3$.

A point of the frame cover carries a frame $v:V\to E_s[12]$. Lift $a$ and
$d$ by retaining this frame under the identity elliptic identification, and
lift $b$ by transporting it to $\phi_{s^3}\circ v$. For $h\in G$, let $h$
act by reparameterizing the frame while leaving elliptic points fixed.
Write $H_-$ for $h=-\mathrm{Id}$ and $C$ for simultaneous inversion of all
four fibre points, with the frame fixed. Upper-case $A,B,D$ denote the
pullback correspondences associated to the lower-case maps; $g^*$ is
represented by the transpose of its graph, with the pullback composition
convention used throughout.

The coordinate formulas give the following identities, already on the open
family with frames:
\[
\begin{gathered}
A^3=1,\qquad B^2=D^2=1,\qquad AD=DA,\qquad BAB=A^{-1},\\
\boxed{BDB=H_-CD.}
\end{gathered}
\]
For the last formula, applying $b,d,b$ sends
$(s,v,P_1,\ldots,P_4)$ to $(-s,-v,-P_1,\ldots,-P_4)$.
Omitting the frame sign or asserting $BD=DB$ on the whole space would be an
error. Both signs will disappear on the precise summand used below.

\textbf{Extension across the cusps.} An isomorphism of the generic elliptic
curve extends to its minimal regular model by uniqueness \cite{StacksMinimal}. This applies
also to the translated sections and to base automorphisms permuting cusps.
The local extensions agree on the dense open set and hence glue. The maps
then preserve the nodal strata used in the canonical fibre-product resolution,
so all act as genuine automorphisms of $X$, and their displayed relations
remain exact there. This uses that specific resolution, not an arbitrary one.

Frame reparameterizations commute with $A,B,D$. They normalize the universal
translation subgroup. The maps $A,B,D$ transport the frame and fibre points
together, so they normalize the translations, individual inversions, and
permutations, preserving the character used in the following construction.

All relevant automorphisms generate a finite group: the translation, inversion,
permutation and frame groups are finite, and the displayed relations give a
finite quotient generated by $A,B,D$. We may therefore use their rational
group algebra and its graph correspondences without a conjectural lifting of
cohomological idempotents.


\subsection{Exact rational Chow projectors}
Let
\[
\Gamma=(V\rtimes\{\pm1\})^4\rtimes\mathfrak S_4.
\]
Its character $\epsilon$ is trivial on translations, is the product of the
four inversion signs, and is $\sgn$ on $\mathfrak S_4$. Define
\[
\Pi_\epsilon=\frac1{|\Gamma|}\sum_{\gamma\in\Gamma}
\epsilon(\gamma)\gamma^*,\qquad
e_G=\frac1{|G|}\sum_{h\in G}h^*,\qquad q=e_G\Pi_\epsilon.
\]
The signed permutation is essential: the Koszul sign from exchanging degree-one
cohomology factors makes this select $\Sym^4 R^1$, rather than its alternating
power. The translation average supplies the component condition at the boundary.

The group $G$ normalizes $\Gamma$ and preserves $\epsilon$. Thus $e_G$ and
$\Pi_\epsilon$ commute, and $q^2=q$. Moreover $q$ commutes with $A,B,D$.
There are exact identities in the correspondence algebra
\[
H_-q=q,\qquad Cq=\epsilon(C)q=(-1)^4q=q.
\]
Hence $BDBq=Dq$, and $BDq=DBq$. Define
\[
d_+=\frac{1+D}{2},\qquad e_\nt=\frac{2-A-A^2}{3},\qquad
b_\pm=\frac{1\pm B}{2}
\]
Their restrictions $qd_+$, $qe_\nt$, and $qb_\pm$ are commuting idempotents
in $q\Q[\mathcal G]q$, whose identity is $q$ and where $\mathcal G$ is the
finite group just constructed. Indeed $B$ conjugates $A$ to $A^{-1}$, so it
commutes with $e_\nt$; $D$ commutes with $A$; and the remaining commutator
has just been removed by $q$. Consequently
\[
\boxed{p_\pm=q\,\frac{1+D}{2}\,\frac{2-A-A^2}{3}\,\frac{1\pm B}{2}}
\]
satisfy $p_\pm^2=p_\pm$, $p_+p_-=p_-p_+=0$, and
$p_++p_-=qd_+e_\nt$. These are equalities modulo rational equivalence,
because they already hold for the finite sums of automorphism correspondences.

\textbf{Descent to $\Q$.} The frame torsor, $X$, $B,D$, and $q$ are defined
over $\Q$. Although $A$ is defined over $\Q(\zeta)$, its Galois conjugate
is $A^2$. Therefore $A+A^2$, $e_\nt$, and $p_\pm$ are rational algebraic
correspondences. More formally, restriction and corestriction of rational Chow
cycles descend these invariant sums; restriction is injective after tensoring
with $\Q$. Thus the idempotent equalities hold over $\Q$, not only over
$\Q(\zeta)$.

No standard conjecture is used here. This is a decomposition by actual finite
automorphisms; it does not assert a Hecke-eigenform Chow decomposition or any
irreducibility of the resulting rank-four realizations by itself.
Section~\ref{sec:irreducibility} gives that separate argument for $\ell\ne3$.


\subsection{Parabolic realization and comparison}
Scholl's finite-group construction \cite[Theorem~1.2.1]{SchollMotives}, in the semistable
setting justified in \cite[Sections~2.7--2.12]{SchollVanishing}, gives for Betti and etale cohomology
\[
\Pi_\epsilon H^j(X)=0\quad(j\ne5),\qquad
\Pi_\epsilon H^5(X)\simeq
H^1_!(S',\Sym^4 R^1\pi'_*\Q).
\]
Here and below the coefficients on the etale side are $\Q_\ell$ and the
base is geometrically extended. The identification is induced by restriction
from the compact variety to the open fibre product, then relative Kunneth
and Leray. It is compatible with every family automorphism above.

There is a useful descent subtlety. Some choices identifying the cusp terms
in Scholl's boundary ladders become equivariant only after a finite extension.
We do not use those chosen boundary identifications. The proper-to-open
restriction, and the assertion that its image is the parabolic image, are
natural. Their injectivity and image equality may be verified after a finite
extension, and then hold for the same maps over $\Q$. The resulting
identification is fully $G_\Q$-equivariant. Nor is it necessary, or correct,
to say that every boundary cohomology group is killed by $\Pi_\epsilon$.

For a finite etale torsor with a finite group, pullback and trace identify
rational cohomology downstairs with invariants upstairs. The statement holds
for both ordinary and compactly supported cohomology, commutes with the
natural map between them, and therefore holds for its image. First take $G$
invariants and then $D$ invariants. This gives
\[
\begin{aligned}
qH^5(X)&\simeq H^1_!(S,\Sym^4 R^1\pi_{S*}\Q),\\
qd_+H^5(X)&\simeq H^1_!(T,\Sym^4 R^1\pi_*\Q).
\end{aligned}
\]
On the latter space $A$ is exactly the manuscript's cubic deck pullback.
The map $B$ uses $u=s^3$ in $\phi_u$, so its induced map is exactly $I$.
The discrepancy under $s\mapsto-s$ is simultaneous elliptic inversion,
already killed on $\Sym^4$. Applying $e_\nt$ and $b_\pm$ proves
\[
p_\pm H^5(X)\simeq(e_\nt H^1_!(T,\Sym^4 R^1\pi_*\Q))^{I=\pm1}.
\]
This proves the theorem's etale assertion for every $\ell$, in particular
the required equality with the existing $2$-adic $W_4^\pm$.

All maps used are pullback, trace, algebraic automorphisms, and natural
cohomological maps. Betti--etale comparison is functorial for these operations,
so the same formula identifies the rational Betti realization. Since it is a
projected subspace of $H^5$ of a smooth projective variety, its Hodge structure
is pure of weight five. Its precise Hodge numbers are computed next.


\subsection{Hodge numbers and an independent sign check}
The Legendre classifying map to the moduli \emph{stack} of elliptic curves is
etale: it factors as a $C_2$-torsor over the full-level-2 stack followed by an
$\mathfrak S_3$-torsor \cite[Lemma~4.3 and Proposition~4.5]{AntieauMeier}. This includes the
elliptic points where the coarse $j$-map ramifies. The maps $s\mapsto s^6$
and $S'\to S$ are etale on the open bases. Apply Scholl \cite[Section~2.13]{SchollVanishing} on $S'$,
where all the section hypotheses were verified, then take $G$ invariants.
The parabolic Hodge types are only $(5,0)$ and $(0,5)$, with
\[
H^{5,0}_!(S,\Sym^4 R^1\pi_{S*})
\simeq H^0(\PP^1_s,\omega^4\otimes\Omega^1_{\PP^1_s}),
\]
where $\omega$ is the Hodge line bundle on $\PP^1_s$. To justify this section
formula after descent, write $r:\overline S'\to\PP^1_s$. Semistable base
change gives $\omega'=r^*\omega$, and
$(r_*\Omega^1_{\overline S'})^G=\Omega^1_{\PP^1_s}$; locally this is the
identity $ds=e u^{e-1}du$ for $s=u^e$ in characteristic zero. The projection
formula then gives the displayed equality of holomorphic spaces.

The singular-fibre multiplicities sum to $12+12+6\cdot2=36$, giving
$\deg\omega=3$ and $\omega^4\otimes\Omega^1\simeq\mathcal O(10)$.
The same degree follows directly from the two minimal charts. Write
$\omega_s=dx/y$ and $\omega_t=dX/Y$ for relative invariant differentials.
On their overlap $\omega_s=t^3\omega_t$. Thus
\[
\beta_j=s^j\,ds\,\omega_s^4
       =-t^{10-j}\,dt\,\omega_t^4,\qquad 0\le j\le10,
\]
are a basis of holomorphic sections; the exponents show regularity at infinity.
The fourth power denotes the symmetric coefficient power, not a fourth wedge
of a one-form on the same elliptic curve.

Since $D^*\beta_j=(-1)^{j+1}\beta_j$, its invariant part has $j=1,3,5,7,9$.
With $z=s^2$ and $\omega_z=dx/y$ on the open family over $T$, let
\[
\alpha_k=z^k\,dz\,\omega_z^4\quad(0\le k\le4);
\qquad (s\mapsto s^2)^*\alpha_k=2\beta_{2k+1}.
\]
These give the five-dimensional $(5,0)$ part on $T$. The literal maps yield
\[
A\alpha_k=\zeta^{k+1}\alpha_k,\qquad I\alpha_k=-\alpha_{4-k}.
\]
For the second identity, use the target coordinate $s'$: $B^*\omega_{s'}=s^3\omega_s$
and $B^*(ds')=d(s^{-1})=-s^{-2}ds$, hence $B^*\beta_j=-\beta_{10-j}$.
This fixes the sign without any appeal to numerical traces.

The projector $e_\nt$ removes precisely $\alpha_2$. On the remaining four
holomorphic dimensions the two eigenspaces are
\[
\begin{array}{c|c}
I=+1&\alpha_0-\alpha_4,\quad\alpha_1-\alpha_3\\[2pt]
I=-1&\alpha_0+\alpha_4,\quad\alpha_1+\alpha_3.
\end{array}
\]
Each has dimension two. Rationality of $p_\pm$ and complex conjugation give
two antiholomorphic dimensions as well. Endpoint Hodge type excludes all
other $(p,q)$. Each realization therefore has rank four and the asserted
Hodge numbers, independently confirming the prior rank calculation.


The construction and the Hodge calculation complete the proof of
Theorem~\ref{thm:scholl-realization}.
